\documentclass[11pt]{article}
\usepackage[left=0.8in,right=0.8in,top=0.9in,bottom=1in]{geometry}
\usepackage{times}
\usepackage{latexsym}
\usepackage[T1]{fontenc}
\usepackage[utf8]{inputenc}
\usepackage{graphicx}

\title{CSE 5525 Final Project}

\author{
\begin{tabular}{ccc}
James Huang & Wilson Wang & Dilmar Roblero \\
\texttt{huang.4978@osu.edu} & \texttt{wang.15722@osu.edu} & \texttt{roblero.7@osu.edu}
\end{tabular}\\[0.5em]
}

\begin{document}
\maketitle

\section{Project Goals}
The goal of this project is to reproduce and analyze a simplified post-training pipeline for large language models consisting of supervised fine-tuning followed by preference-based alignment using Direct Preference Optimization. We will then analyze where and why the alignment helps/hurts, and what tradeoffs it introduces between helpfulness and safety, and how design decisions such as chat template choice and data quality can affect the model's capabilities.

\section{NLP Task}
The task of this project is to generate a helpful, accurate, and safe response $y$ for a given instruction or user prompt $x$. This project will be evaluated on four evaluation sets: GSM8K (mathematical reasoning), MBPP (code generation), IFEval (instruction following), and HarmBench (safety). 

\section{Data}
\paragraph{Instruction data (SFT):}
\begin{itemize}
  \item Dataset: Tulu 3/OLMo 2 SFT
  \item Contains 866,138 examples
\end{itemize}

\paragraph{Preference data:}
\begin{itemize}
  \item Dataset: OLMo 2 1B preference dataset
  \item Contains 378,301 examples
\end{itemize}

\section{Methods}
All training is performed on the base model (Llama-3.2-1B) via the Tinker API.
We also evaluate (Llama-3.2-1B-Instruct) as a reference model to provide a good direction on how an industry-trained and instruction-tuned model developed by Meta would perform.
\paragraph{Part A: Supervised Fine-Tuning (SFT)}
We fine-tune the base model on the Tulu 3 SFT dataset. Our baseline configuration uses: learning rate 0.0005, linear learning-rate schedule, 1 epoch, batch size 128, max sequence length 16384, and LoRA rank 64. We train with default render\_name (\texttt{role\_colon} chat template), which renders conversations using plain-text role labeling (\texttt{User:} / \texttt{Assistant:}). 
We evaluate two checkpoints from this run: one at 2500 steps and one at the end of the full epoch. This can help us understand how training on the full dataset is important for the model's capabilities. Additionally, we train a separate full dataset variant using the Llama3 chat template (with special tokens like \texttt{<|start\_header\_id|>}) and learning rate 0.0001 to see the effect of template choice on model performance.
\paragraph{Part B: Alignment with Preferences}
We will start from our best SFT checkpoint, and run DPO using the OLMo~2 preference dataset. DPO directly optimizes the model so that preferred responses receive higher probability than dispreferred ones. The key hyperparameter is $\beta$, the Kullback-Leibler (KL) penalty coefficient, which controls how far the aligned policy can deviate from the SFT reference policy. We will begin with $\beta = 0.1$ as our baseline, and then compare SFT vs.\ SFT+DPO on all four evaluation sets, focusing on tradeoffs in instruction-following, refusal behavior, factuality, and robustness.
 
\section{Extensions}
\begin{itemize}
    \item \textbf{Safety vs.~helpfulness tradeoff with error analysis.}
    We will sweep across $\beta \in \{0.01, 0.05, 0.1, 0.3, 0.5\}$ and plot refusal rate vs.~task performance. Beyond that, we will manually categorize a sample of outputs from the four evaluation sets into failure types (e.g., hallucination, over-refusal, and verbosity collapse) and compare how the failure distribution shifts before alignment (SFT) and after alignment (best DPO).

    \item \textbf{Data quality filtering.}
    We will first compute quality signals (length, duplicates, toxicity) on the SFT dataset, then build two training sets: (a) the original examples dataset and (b) a filtered dataset with examples meeting all quality filters, such as removing short/long responses, duplicated instruction-response pairs, and toxic responses. We will then measure how data quality affects both SFT and downstream DPO performance.

    \item \textbf{LoRA rank design analysis (if budget permits).}
    We will train with different LoRA rank values and compare evaluation scores, training loss curves, and credit cost to find the best balance between model quality and training cost.
\end{itemize}

\section{Baselines}
\begin{itemize}
  \item Base model: Llama-3.2-1B without any fine-tuning.
  \item SFT model: The 2500-step and the full-dataset SFT models trained on the Tulu 3 SFT dataset with \texttt{role\_colon} chat template, plus one full-dataset variant using the Llama3 template.
  \item SFT + DPO model: The SFT model with preference alignment on the OLMo 2 1B preference dataset.
\end{itemize}

\section{Evaluation}
We will use the OLEMS evaluation system, which runs over the following evaluation sets: 
\begin{itemize}
  \item GSM8K: grade school math word problems (mathematical reasoning)
  \item MBPP: mostly Basic Python Problems (code generation)
  \item IFEval: instruction following evaluation
  \item HarmBench: safety and harmfulness evaluation
\end{itemize}

\section{Check-in Results}
\subsection{SFT Results}
Table~\ref{tab:sft_results} shows our SFT results on four benchmarks. Compared with the untuned baseline, both SFT checkpoints improve performance on GSM8K, MBPP, and IFEval, confirming that supervised fine-tuning helps the model better follow instructions and produce stronger task responses. We include Llama-3.2-1B-Instruct as a reference model because it 
provides a strong instruction-tuned comparison point. In addition, we report one full-dataset SFT variant using the Llama3 chat template as an additional comparison.
\begin{table}[htbp]
\centering
\small
\begin{tabular}{lccccc}
\hline
\textbf{Benchmark} & \textbf{Baseline} & \textbf{SFT 2500-step} & \textbf{SFT Full} & \textbf{SFT Full +} & \textbf{Llama-Instruct} \\
 &  &  &  & \textbf{Llama3} &  \\
\hline
GSM8K     & 3.6\%  & 9.5\%  & 11.7\% & 6.3\%  & 42.0\% \\
MBPP      & 20.8\% & 22.4\% & 24.0\% & 21.6\% & 32.4\% \\
IFEval    & 18.3\% & 40.9\% & 46.4\% & 43.8\% & 65.2\% \\
HarmBench & --     & 84.7\% & 82.2\% & 64.1\% & 83.4\% \\
\hline
\end{tabular}
\caption{SFT results on four benchmarks. Llama-3.2-1B-Instruct is included as a reference model.}
\label{tab:sft_results}
\end{table}

\subsection{Analysis and Experiments}
\paragraph{Why the Llama~3 chat template underperformed.}
We used the Llama 3 chat template because it is the same format adopted by Llama-3.2-1B-Instruct, which is fine-tuned by Meta. However, the full-dataset variant using the Llama3 template performs worse than the \texttt{role\_colon} variant on all benchmarks, especially on GSM8K and HarmBench. A likely reason is that the Llama 3 template relies on special tokens (\texttt{<|start\_header\_id|>}, \texttt{<|eot\_id|>}) that exist in the vocabulary but were never meaningfully trained during pretraining of the base Llama-3.2-1B model. As a result, during SFT the model must simultaneously learn the semantics of these unfamiliar tokens and learn to follow instructions. This creates a more difficult task compared to using a plain-text template like \texttt{role\_colon}, where the role labeling (\texttt{User:} / \texttt{Assistant:}) are already well-represented in the pretraining corpus.

\subsection{Training Cost}
So far, our total training cost is approximately \$153 with \$88 spent on SFT + Llama3 chat template and \$65 on SFT + \texttt{role\_colon} template. We still have \$597 remaining in our team budget.

\subsection{Comparison and Model Selection}
We compared the 2500-step, full-dataset, and Llama3 chat-template SFT variants in Table~\ref{tab:sft_results}. The full-dataset SFT model performs better than the 2500-step checkpoint on GSM8K, MBPP, and IFEval, and also outperforms the additional Llama3 chat template variant overall. Based on this comparison, we select the full-dataset SFT model with \texttt{role\_colon} chat template as our best SFT model for the next DPO stage. 


\subsection{Future Plan}

\begin{itemize}
    \item \textbf{By April 5:} Complete the DPO stage and prepare the aligned model for further evaluation.
    
    \item \textbf{April 5--April 13:} Work on two extensions in parallel:
    \begin{itemize}
        \item \textbf{Safety vs.~helpfulness tradeoff with error analysis} 
        \item \textbf{Data quality filtering} 
    \end{itemize}
    
    \item \textbf{By April 13:} Complete the main experiments for both extensions.
    
    \item \textbf{April 13--April 20:} Finish extension analysis, organize results, and complete the presentation slides and final report.
    
    \item \textbf{By April 20:} Submit the final project report and complete the last extension.
\end{itemize}

\section{Ethical Challenges}
\paragraph{Risk of harmful content generation and misuse.}
The safety evaluation prompts may contain harmful and sensitive content. During generation and evaluation, a misaligned model could produce responses that are actionable and could be misused. To mitigate this risk, we will keep the model, evaluation prompts, and raw generations private to the team, Professor and TA. 

\paragraph{Implicit bias in preference-based alignment.}
Preference optimization fixes an implicit definition of what counts as a ``good'' response, but this definition may reflect cultural preferences or biases from the dataset creators. As a result, DPO may amplify bias, making the model more favorable toward certain viewpoints or less supportive toward certain forms of expression. As we mentioned above, We will not deploy the model and will use it strictly for this course project. In our analysis, we will examine outputs for potential bias patterns and document representative examples.

\paragraph{Overconfidence and factuality regressions.}
Alignment may increase the model's fluency and apparent confidence without improving factual correctness, potentially exacerbating hallucinations and misleading users. As a mitigation, we will encourage calibrated responses for uncertain queries (e.g. expressing uncertainty, requesting clarification, or recommending verification).

\end{document}
